\documentclass[10pt,twocolumn]{article}
\usepackage[T1]{fontenc}
\usepackage[utf8]{inputenc}
\usepackage{lmodern}
\usepackage[margin=1.8cm,top=2.4cm,bottom=2cm,columnsep=0.6cm]{geometry}
\usepackage{fancyhdr}
\usepackage{titlesec}
\usepackage{booktabs}
\usepackage{amsmath,amssymb,amsfonts}
\usepackage{microtype}
\usepackage{xcolor}
\usepackage{mdframed}
\usepackage{parskip}
\usepackage{enumitem}
\usepackage{hyperref}

\definecolor{rulered}{RGB}{180,30,30}
\definecolor{boxgray}{RGB}{245,245,245}
\definecolor{keywordgray}{RGB}{100,100,100}

\pagestyle{fancy}
\fancyhf{}
\fancyhead[L]{\textsc{\small Claude Mag}}
\fancyhead[C]{\small\textit{Machine Learning Research Digest}}
\fancyhead[R]{\small 2026-09-16}
\fancyfoot[C]{\small\thepage}
\renewcommand{\headrulewidth}{0.8pt}

\titleformat{\section}{\normalsize\bfseries\color{rulered}}{}{0pt}{}%
  [\vspace{-4pt}\color{rulered}\rule{\columnwidth}{0.4pt}]
\titlespacing{\section}{0pt}{8pt}{4pt}

\hypersetup{colorlinks=true,urlcolor=rulered,linkcolor=black}

\begin{document}
\setlength{\parindent}{0pt}
\setlength{\parskip}{4pt}

{\small\color{keywordgray}\textbf{Keywords:}
  motion generation \textbullet{}
  point cloud representation \textbullet{}
  text-to-motion \textbullet{}
  animal motion \textbullet{}
  diffusion models}\par\vspace{4pt}

{\LARGE\bfseries\raggedright
  UniMo: Unifying Human and Animal
  Motion Generation}\par\vspace{4pt}

{\large\itshape\raggedright
  Parametric Skeletons Block Cross-Species Unification ---
  Topology-Agnostic Point Clouds and a Dataset 102$\times$ Larger
  Than Any Prior Animal Corpus Enable a Single Model for All
  Species}\par\vspace{6pt}

{\small\textbf{By} Zeyu Zhang, Zhiyuan Zhang, Siheng Wang, Yiran Wang,
  Danning Li, Ian Reid, Richard Hartley
  (University of Adelaide; Australian National University)}\par\vspace{2pt}

{\small\color{keywordgray}arXiv:2609.12342 --- SIGGRAPH Asia 2026}
\par\vspace{2pt}\noindent{\color{rulered}\rule{\columnwidth}{0.6pt}}\vspace{6pt}

Text-driven 3D motion generation has matured rapidly for humans, with
standardized SMPL skeletons enabling a community of comparable models.
Animals remain largely excluded: a dog has 17 joints, a horse 30, a
bird 15 --- their skeletal topologies are incompatible with each other
and with SMPL, forcing per-species model training. UniMo eliminates
this barrier by converting any skeleton to a fixed-size unparametric
point cloud, training a single diffusion model across species, and
releasing UniML3D, a dataset that is 102 times larger than any
existing animal motion corpus.

\begin{mdframed}[backgroundcolor=boxgray,linecolor=rulered,linewidth=1pt,
    innerleftmargin=6pt,innerrightmargin=6pt,
    innertopmargin=6pt,innerbottommargin=6pt]
\textbf{\small AT A GLANCE}\par\vspace{4pt}
\begin{description}[leftmargin=0pt,labelsep=4pt,itemsep=2pt,parsep=0pt,
    font=\small\bfseries]
  \item[Problem:] Diverse animal skeletal topologies prevent training
    a single motion generation model across species.
  \item[Why it matters:] Per-species models are expensive; animal
    motion datasets are tiny ($<$1,500 sequences); SIGGRAPH
    applications demand diverse creatures.
  \item[Method:] Convert parametric skeletons to $N=256$ point clouds
    with dynamic sampling; train single diffusion model; IK
    reconstruction per species at test time.
  \item[Dataset:] UniML3D: 145,907 sequences + 433,388 captions
    across human and animal categories.
  \item[Result:] State of the art on HumanML3D, KIT-ML, and
    AnimalML3D simultaneously in one model.
\end{description}
\end{mdframed}

\section{Why Existing Approaches Fall Short}

Motion generation models represent poses as joint rotation/position
vectors in species-specific fixed order:
\[
  \mathbf{p}_\text{human} \in \mathbb{R}^{22 \times 3}, \quad
  \mathbf{p}_\text{dog} \in \mathbb{R}^{17 \times 3}, \quad
  \mathbf{p}_\text{horse} \in \mathbb{R}^{30 \times 3}
\]
Dimension and joint semantics differ across species; concatenating
these into a single input space is undefined without topology alignment
that does not exist. The result: every new species requires a separate
model trained from scratch.

\section{Core Mechanism: Point Cloud Representation}

UniMo samples $N = 256$ points from the bones of any skeleton,
regardless of species topology. For each bone $b$ connecting joints
$j_a$ and $j_b$, $n_b$ equally-spaced points are placed along the
bone segment. The pose at time $t$ is:
\[
  \mathbf{P}_t \in \mathbb{R}^{N \times 3}, \quad N = \sum_b n_b
\]

This representation is topology-agnostic (always 256 points),
continuous (captures limb geometry, not just endpoints), and
differentiable (linear interpolation along bones).

\textbf{Dynamic sampling.} Bone activity at time $t$:
\[
  v_{b,t} = \|\Delta\mathbf{p}_{j_a,t}\|_2 + \|\Delta\mathbf{p}_{j_b,t}\|_2
\]
Mean activity: $\bar{v}_b = \frac{1}{T}\sum_t v_{b,t}$.
Point allocation:
\[
  n_b = \max\!\left(1,\;\left\lfloor N\cdot
    \frac{\bar{v}_b}{\sum_{b'}\bar{v}_{b'}}\right\rfloor\right)
\]
This concentrates the $N$-point budget on actively moving bones,
sharpening the representation of the most informative parts of motion.

\section{Technical Formulation}

\textbf{Diffusion forward process:}
\[
  \mathbf{P}_t^{(\tau)} = \sqrt{\bar{\alpha}_\tau}\,\mathbf{P}_t^{(0)}
    + \sqrt{1-\bar{\alpha}_\tau}\,\boldsymbol{\epsilon},
    \quad\boldsymbol{\epsilon}\sim\mathcal{N}(0,I)
\]

\textbf{Denoising objective:}
\[
  \mathcal{L} = \mathbb{E}_{\tau,\boldsymbol{\epsilon}}\!\left[
    \left\|\boldsymbol{\epsilon}
      - \boldsymbol{\epsilon}_\theta\!\left(\mathbf{P}^{(\tau)},
        \tau, \mathbf{c}\right)
    \right\|^2\right]
\]
where $\mathbf{c}$ is a CLIP text embedding and
$\boldsymbol{\epsilon}_\theta$ is a Transformer denoiser with
cross-attention to $\mathbf{c}$.

\textbf{Inverse kinematics reconstruction:} Given predicted point cloud
$\hat{\mathbf{P}}_t$, recover joint rotations by minimizing:
\[
  \hat{\boldsymbol{\theta}}_t = \arg\min_{\boldsymbol{\theta}}
    \sum_b\!\left\|\mathrm{FK}(\boldsymbol{\theta})_b
      - \hat{\mathbf{P}}_{t,b}\right\|^2
    + \lambda\sum_b\!\left(\|\Delta\mathrm{FK}_b\|_2 - \ell_b\right)^2
\]
where $\mathrm{FK}$ is forward kinematics and $\ell_b$ is rest-pose
bone length.

\section{UniML3D Dataset}

\begin{center}
\begin{tabular}{lrr}
\toprule
Category & Sequences & Captions \\
\midrule
Human             & 128,319 & 384,957 \\
Animal (quadruped) & 14,231 & 42,693 \\
Animal (avian)     & 3,357  & 5,738  \\
\midrule
\textbf{Total}    & \textbf{145,907} & \textbf{433,388} \\
\bottomrule
\end{tabular}
\end{center}

The largest prior animal dataset contains $\sim$1,432 sequences;
UniML3D is 102$\times$ larger.

\section{Experimental Results}

\begin{center}
\begin{tabular}{lccc}
\toprule
Method & HumanML3D & KIT-ML & AnimalML3D \\
       & FID$\downarrow$ & FID$\downarrow$ & FID$\downarrow$ \\
\midrule
MDM            & 0.544 & 0.497 & --- \\
MotionDiffuse  & 0.630 & 1.954 & --- \\
MLD            & 0.473 & 0.404 & --- \\
Species-specific & ---  & ---  & 2.14 \\
\textbf{UniMo} & \textbf{0.412} & \textbf{0.381} & \textbf{1.07} \\
\bottomrule
\end{tabular}
\end{center}

UniMo improves over MLD on both human benchmarks while being the
first unified model to achieve competitive quality on animals.

\section{Ablations and Interpretation}

\textbf{Without dynamic sampling (static uniform):} AnimalML3D FID
degrades by 0.11; most degradation on fast-moving quadruped sequences.

\textbf{Without joint training on animals:} AnimalML3D FID $= 4.32$
(near random), confirming human-only training does not generalize.

\textbf{Point count $N$:} Plateaus at $N=256$; $N=128$ degrades FID
by 0.09 on animals; $N=512$ shows no improvement.

\textbf{IK reconstruction:} Mean bone-length error of 1.2~mm
($<$0.1\% of body height), confirming the IK solver accurately
recovers valid joint configurations from predicted point clouds.

\section*{Reference}

Zeyu Zhang et al.\ \textbf{UniMo: Unifying Human and Animal Motion
Generation.} arXiv:2609.12342, September 2026. SIGGRAPH Asia 2026.
\url{https://arxiv.org/abs/2609.12342}

\end{document}
